\documentclass{article}
\usepackage[UTF8]{ctex}
\usepackage{amsmath,amssymb,amsfonts}
\usepackage{enumitem}
\usepackage{geometry}
\geometry{a4paper, margin=1in}
% === 错题本排版优化（不参与题目标识） ===
\setlength{\parindent}{0pt}
\setlength{\parskip}{0.35em plus 0.10em minus 0.08em}
\setlength{\emergencystretch}{3em}
\setlist{itemsep=0.25em, topsep=0.40em, parsep=0pt}
\allowdisplaybreaks[2]
% === 错题本排版优化结束 ===
\begin{document}
\section*{浙江大学2024年高等代数真题}
\section*{94．浙江大学2024年高等代数试卷}

\subsection{题目1}
1．假设 \(\frac{5}{11}\) 是整系数多项式 \(f(x)\) 的根，证明：\(f(\sqrt{-1}) f(-\sqrt{-1})\) 是正整数，且是146的倍数。

\subsection{题目2}
2．求线性方程组

\[
\left\{\begin{array}{r}
x_{1}-2 x_{2}+3 x_{3}-x_{4}=0 \\
3 x_{1}+x_{2}+x_{3}-3 x_{4}=0 \\
x_{1}+5 x_{2}-5 x_{3}-x_{4}=0
\end{array}\right.
\]

的基础解系，假设该方程组的解和另外一个解为 \(k_{1}\left(\begin{array}{c}3 \\ -4 \\ -3 \\ 0\end{array}\right)+k_{2}\left(\begin{array}{l}0 \\ 2 \\ 2 \\ 1\end{array}\right)\) 的方程组有公共解，求出所有公共解。

\subsection{题目3}
3．设 \(X=\left(x_{1}, x_{2}, x_{3}, x_{4}\right), X^{\mathrm{T}}\) 是 \(X\) 的转置，问是否存在一次多项式

\[
u_{i}(x)=a_{i} x_{1}+b_{i} x_{2}+c_{i} x_{3}+d_{i} x_{4}(i=1,2,3,4)
\]

满足

\[
X\left(\begin{array}{llll}
2 & 0 & 2 & 4 \\
0 & 2 & 4 & 2 \\
2 & 4 & 2 & 0 \\
4 & 2 & 0 & 2
\end{array}\right) X=\left|\begin{array}{ll}
u_{1}(x) & u_{2}(x) \\
u_{3}(x) & u_{4}(x)
\end{array}\right| .
\]

并说明理由．

\subsection{题目4}
4．假设 \(\mathbb{R}\) 是实数域，实向量空间 \(\mathbb{R}^{3}\) 两组向量分别为 \(\alpha_{1}=(-1,1,0), \alpha_{2}=(2,-1,2), \alpha_{3}=(0,1, b)\) 和

\[
\beta_{1}=(1,0,-1), \beta_{2}=(-1,1,1), \beta_{3}=(1,1, c) .
\]

（1）当 \(b, c\) 取何值时，不存在 \(\mathbb{R}^{3}\) 上的线性变换 \(\mathcal{T}\) ，满足 \(\mathcal{T}\left(\alpha_{i}\right)=\beta_{i}, i=1,2,3\) 。\\
（2）当 \(b, c\) 取何值时，至少存在两个 \(\mathbb{R}^{3}\) 上的线性变换 \(\mathcal{T}\) 满足 \(\mathcal{T}\left(\alpha_{i}\right)=\beta_{i}, i=1,2,3\) ．\\
（3）当 \(b, c\) 取何值时，存在 \(\mathbb{R}^{3}\) 上唯一的线性变换 \(\mathcal{T}\) ，满足 \(\mathcal{T}\left(\alpha_{i}\right)=\beta_{i}, i=1,2,3\) ．这样的线性变换是正交变换吗？为什么？

\subsection{题目5}
5．假设 \(A\) 是2024阶方阵，主对角线上全是偶数，其余的都是奇数。证明该矩阵为可逆矩阵。

\subsection{题目6}
6．求 \(f(x)=x^{4}+2 x^{3}-x^{2}-4 x-2, g(x)=x^{4}+x^{3}-x^{2}-2 x-2\) 的最大公因式 \(d(x)\) 以及多项式\(u(x), v(x)\) 满足 \(d(x)=u(x) f(x)+v(x) g(x)\) ．

\subsection{题目7}
7．已知矩阵

\[
A=\left(\begin{array}{ccc}
3 & a & 6 \\
-5 & -7 & -1 \\
b & -6 & -7
\end{array}\right)
\]

的特征值 \(\lambda\) 对应的特征向量 \(\alpha=\left(\begin{array}{c}-2 \\ 1 \\ 2\end{array}\right)\) ，求该矩阵的若当（Jordan）标准型．

\subsection{题目8}
8．设 \(A\) 是 \(n\) 阶非幂零矩阵，\(S\) 是使得 \(\lambda A\) 与 \(\lambda\) 相似的复数 \(\lambda\) 的集合，证明 \(S\) 是一个有限集．

\subsection{题目9}
9．设 \(\alpha, \beta, \gamma\) 是有理数域上线性空间 \(V\) 中的向量，其中 \(\alpha \neq 0\) ，假如存在 \(V\) 上线性变换 \(\mathcal{T}\) ，使得

\[
\mathcal{T} \alpha=\beta, \mathcal{T} \beta=-\gamma, \mathcal{T} \gamma=\alpha-\beta .
\]

证明：\(\alpha, \beta, \gamma\) 在 \(V\) 中线性无关．

\subsection{题目10}
10．已知 \(\mathcal{A}, \mathcal{B}, \mathcal{C}\) 是有限维线性空间上的三个线性变换，证明： \(\mathcal{A}+\mathcal{B}-\mathcal{A} \mathcal{C} \mathcal{B}\) 和 \(\mathcal{A}+\mathcal{B}-\mathcal{B} \mathcal{C} \mathcal{A}\) 核空间是同构的．
\end{document}
